
\medskip

%Une entreprise fabrique des objets en plastique en injectant dans un moule de la matière fondue à $210$~\degres C. On cherche à modéliser le refroidissement du matériau à l'aide d'une fonction $f$ donnant la température du matériau injecté en fonction du temps $t$.
%
%Le temps est exprimé en seconde et la température est exprimée en degré Celsius.
%
%On admet que la fonction $f$ cherchée est solution d'une équation différentielle de la forme suivante où $m$ est une constante réelle que l'on cherche à déterminer:

\[(E) : \quad  y'  +0,02y = m\]

\medskip

\textbf{Partie A}

\medskip

\begin{enumerate}
\item %Justifier l'affichage suivant d'un logiciel de calcul formel:

%\begin{center}
%\begin{tabularx}{\linewidth}{|l|X|}\hline
%Entrée :& RésoudreEquationDifférentielle $(y' + 0,02y = m)$\\ \hline
%Sortie :& \fbox{$\to$}\: $y = k *\text{exp}(-0.02 * t) + 50 * m$\\ \hline
%\end{tabularx}
%\end{center}
On sait que l'équation différentielle $y' + 0,02y = 0$ a pour solutions les fonctions $t \longmapsto y = k\e^{-0,02t}$, avec $k \in \R$.

D'autre part une fonction constante $y = \alpha$, \:$\alpha \in \R$ est solution de $(E)$ (avec donc $y' = 0$) si $0 + 0,02\alpha = m \iff 0,02\alpha = m \iff \alpha = \dfrac{m}{0,02} = 50m$.

Conclusion : toutes les solutions de $(E)$ sont les fonctions définies par 
\begin{center} $t \longmapsto y = k\e^{0,02t} + 50m$, \: avec $k \in \R$\end{center}
\item %La température de l'atelier est de 30~\degres C. On admet que la température $f(t)$ tend vers 30~\degres C lorsque $t$ tend vers l'infini.

%Démontrer que $m = 0,6$.
On sait que $\displaystyle\lim_{t \to + \infty} \e^{-0,02t} = 0$, donc $\displaystyle\lim_{t \to + \infty} k\e^{-0,02t} = 0$ et $\displaystyle\lim_{t \to + \infty} k\e^{-0,02t} + 50m = 50m$.

Or on a $\displaystyle\lim_{t \to + \infty} k\e^{-0,02t} + 50m  = 30$, donc $m = \dfrac35 = \dfrac{6}{10} = 0,6$.
\item %Déterminer l'expression de la fonction $f$ cherchée en tenant compte de la condition
%initiale $f(0) = 210.$
Avec $50m = 50 \times 0,6 v= 5 \times 6 = 30$, les solutions de $(E)$ sont les fonctions $f$ définies par : $ t \longmapsto f(t) = k\e^{-0,02t} + 30$

On a $f(0) = 210 \iff k \times 1 + 30 = 210 \iff k = 180$.

Finalement : $f(t) = 180\e^{-0,02t} + 30$.
\end{enumerate}

\medskip

\textbf{Partie B}

\medskip

%On admet ici que la température (exprimée en degré Celsius) du matériau injecté en fonction du temps (exprimé en seconde) est donnée par la fonction dont l'expression et une représentation graphique sont données ci-dessous:

\[f(t) =180 \e^{-0,02t} + 30.\]

\begin{center}
\psset{unit=0.05cm,yunit=0.02cm,arrowsize =2pt 3}
\begin{pspicture}(-20,-15)(220,230)
\multido{\n=0+20}{11}{\psline[linewidth=0.2pt](\n,0)(\n,230)}
\multido{\n=0+50}{5}{\psline[linewidth=0.2pt](0,\n)(220,\n)}
\psaxes[linewidth=1.25pt,labelFontSize=\scriptstyle,Dx=20,Dy=50]{->}(0,0)(0,0)(220,230)
\psplot[plotpoints=2000,linewidth=1.25pt,linecolor=blue]{0}{220}{180 2.71828 0.02 x mul exp div 30 add}
\uput[u](190,0){temps (en s)}\uput[r](0,220){température (en \degres C)}
\psline[linestyle=dashed,ArrowInside=->,linewidth=1.25pt](0,50)(109.86,50)(109.86,0)
\psframe[linecolor=red](100,107.82)
\end{pspicture}
\end{center}

\begin{enumerate}
\item %L'objet peut être démoulé lorsque sa température devient inférieure à 50\degres C. 
	\begin{enumerate}
		\item %Par lecture graphique, donner une valeur approchée du nombre $T$ de
%secondes à attendre avant de démouler l'objet.
On lit sur le graphique $T \approx 110$~(s).
		\item %Déterminer par le calcul la valeur exacte de ce temps $T$.
On a $f(T) = 50 \iff 180 \e^{-0,02T} + 30 = 50 \iff 180 \e^{-0,02T} = 20 \iff 20 \times, 9\e^{-0,02T} = 20 \times 1 \iff 9\e^{-0,02T} = 1 \iff \e^{-0,02T} = \dfrac19 \iff -0,02T = \ln \left(\dfrac 19 \right) \iff -0,02T  = - \ln 9 \iff 0,02T = 9 \iff T = \dfrac{\ln 9}{0,02} \approx 109,86$~(s).

Donc $T = \dfrac{\ln 9}{0,02} = 50\ln 9$.
	\end{enumerate}	
\item %À l'aide d'une intégrale, calculer la valeur moyenne de la température sur les $100$~premières secondes.
Soit $\overline{t}$ la valeur moyenne de la température sur les $100$~premières secondes :

$\overline{t} = \dfrac{1}{100}\displaystyle\int_0^{100} \left(180 \e^{-0,02t} + 30\right)\:\text{d}t = \dfrac{1}{100}\left[- \dfrac{180}{0,02}\e^{-0,02t} + 30t \right]_0^{100} = \dfrac{1}{100}\left[- \np{9000}\e^{-0,02t} + 30t \right]_0^{100} = \dfrac{1}{100}\left[-\np{9000}\e^{-2} + \np{3000} - (-\np{9000})\right] = \dfrac{1}{100}\left[-\np{9000}\e^{-2} + \np{12000}\right] = 120 - 90\e^{-2}$ 

soit environ 107,8~(\degres C).

Graphiquement l'aire de la surface limitée par la représentation graphique de $f$, l'axe des abscisses et les droites d'équation $x = 0$ et $x = 100$ est égale à celle du rectangle (en rouge sur la figure) de dimensions 100 et 107,8.
\end{enumerate}

\bigskip

